\section{选择性身份重连}

跟踪器丢了一个人又找回来时，必须判断这个回来的人是不是原来那个。接受，身份继续；拒绝，这个检测重新进入关联池，此后每一个提案都随之改变。我们在这一提交动作之前放一个标量门控。它只读检测器的输出和轨迹状态，不读别的东西：没有外观网络，没有特征库，也不用未来帧。正是这条限制，使方法可以运行在只有检测器和几何跟踪器的场合。

令 $x$ 记录检测置信度、预测框残差、轨迹时长和观测到的间断，以及提出重连那一帧上竞争者一对一分配的代价。记 $m_r$ 与 $m_c$ 为分配矩阵中同一行、同一列里最优竞争者代价减去被选代价，并用 $m=\min(m_r,m_c)$ 表示较弱的一侧。先做基础变换 $z(x)$，再构造交互向量，坐标标准化后拟合
\begin{align*}
g(x) &= b+w^\top \phi(x),\\
\phi(x) &= [z(x),\,m,\,|m_r-m_c|,\,h d,\,h o,\,h m,\,q m,\ldots].
\end{align*}
其中 $h$ 为变换后的间断长度，$d$ 与 $o$ 为归一化的中心位移与框重叠，$q$ 为检测置信度。乘积项让运动一致性与分配歧义各自的权重随间断长度移动，这是纯加性评分无法表达的。$g$ 用在开发集标注提案上的带正则逻辑损失拟合。消融版本另加同源内排序损失
$\log(1+\exp(g_{\rm incorrect}-g_{\rm correct}))$。阈值决定一次重连是否发生；被拒绝的检测以未匹配状态交回跟踪器，下一步关联看到的就是这个决定造成的状态。

\begin{table}[t]
\centering
\footnotesize
\setlength{\tabcolsep}{4pt}
\caption{门控的输入坐标。第 1--12 项构成 $z(x)$，三个竞争摘要把它扩到 15 维，八个乘积项扩到 23 维。剑号表示标准化之前先做 $\log(1+\cdot)$ 变换。$m_r$、$m_c$ 为行、列裕量，$n_r$、$n_c$ 为竞争集合的大小，$h$ 为变换后的间断，$d$、$a$、$o$ 为归一化的中心位移、面积变化与重叠，$q$ 为检测置信度，$q_0$ 为上一次的置信度。}
\label{tab:features-zh}
\begin{tabular}{rll}
\hline
编号 & 特征 & 含义 \\
\hline
1 & \texttt{gap}$^\dagger$ & 轨迹丢失后经过的帧数 \\
2 & \texttt{track\_age}$^\dagger$ & 丢失时轨迹的时长 \\
3 & \texttt{det\_score} & 回归检测的置信度 \\
4 & \texttt{prev\_score} & 上一次接受框的置信度 \\
5 & \texttt{iou} & 与预测框的重叠 \\
6 & \texttt{center\_dist}$^\dagger$ & 归一化中心偏移 \\
7 & \texttt{area\_ratio}$^\dagger$ & 相对上一个框的对数面积比 \\
8 & \texttt{assign\_cost} & 提案对的原始关联代价 \\
9 & \texttt{row\_margin} & 同行最优竞争代价减被选代价 \\
10 & \texttt{col\_margin} & 同列最优竞争代价减被选代价 \\
11 & \texttt{row\_alts}$^\dagger$ & 该行的竞争条目数 \\
12 & \texttt{col\_alts}$^\dagger$ & 该列的竞争条目数 \\
\hline
13 & $\min(m_r,m_c)$ & 竞争中较弱的一侧 \\
14 & $|m_r-m_c|$ & 两侧的不平衡程度 \\
15 & $\max(n_r,n_c)$ & 较大的竞争集合 \\
16 & $h\,d$ & 位移歧义乘以间断 \\
17 & $h\,a$ & 面积变化乘以间断 \\
18 & $h\,o$ & 重叠乘以间断 \\
19 & $h\,m$ & 分配歧义乘以间断 \\
20 & $q\,m$ & 置信度乘以歧义 \\
21 & $q_0\,q$ & 两次置信度的一致程度 \\
22 & $d\,m$ & 位移乘以歧义 \\
23 & $a\,m$ & 轨迹时长乘以歧义 \\
\hline
\end{tabular}
\end{table}

每个坐标都在提交分配之前就可以取得。评分不看未来帧、不看其它来源、也不看任何标注；训练标签是在跟踪器产出提案之后才附加上去的，运行时从不进入门控。

\begin{figure}[t]
\centering
\begin{tikzpicture}[
  node distance=5mm and 8mm,
  box/.style={draw, rounded corners=2pt, align=center, font=\scriptsize,
              inner sep=3.5pt, minimum height=6.5mm},
  arr/.style={-latex, semithick}
]
\node[box] (det) {检测器};
\node[box, right=of det] (asc) {关联级联};
\node[box, right=of asc, fill=blue!10] (gate) {门控\\$g(x)\!\geq\!\tau$};
\node[box, above left=7mm and -2mm of gate] (com) {提交};
\node[box, below left=7mm and -2mm of gate] (rej) {未匹配池};
\draw[arr] (det) -- (asc);
\draw[arr] (asc) -- node[above, font=\tiny] {提案} (gate);
\draw[arr] (gate) -- node[above, font=\tiny] {接受} (com);
\draw[arr] (gate) -- node[below, font=\tiny] {拒绝} (rej);
\draw[arr] (rej.west) -- ++(-3mm,0) |- ([yshift=-3mm]asc.south);
\end{tikzpicture}
\caption{门控的位置。跟踪器在自己的关联级联中提出一次重连；门控只对这一条提案打分，要么提交，要么把检测退回未匹配池。被退回的检测会重新进入下一步关联，因此这个决定改变了后续帧所面对的状​态。}
\label{fig:gate-zh}
\end{figure}

\subsection{如何读这些坐标}

23 个坐标分三组，这个分组也正是消融要分开的东西。前 12 个描述提案本身：轨迹已经消失多久、消失时它有多长、回归检测有多自信、上一次被接受的检测有多自信、预测框贴合得如何、原始关联代价是多少。这些正是手写阈值会用的量，在它们上面的线性评分就是加性基线。

中间 3 个描述的是竞争，而不是提案。一条候选单独看可能很好，却仍然是错误的选择，因为分配矩阵的同一行里坐着更好的候选；$m_r$ 与 $m_c$ 度量这条提案在两侧各有多少余量，而最小值、两侧差值与竞争集合大小概括了这份余量。拥挤画面里的一条提案与空旷画面里一模一样的一条提案，前 12 个坐标相同，后 11 个不同。

最后 8 个是乘积。一个乘积是在说：某项证据在另一项证据所命名的条件下更重要——运动一致性在一帧间断上无关紧要，在五帧间断上却举足轻重；分配歧义在检测自信时比在检测勉强时更值得考虑。加性评分只能把这些效应平均成一个权重；加入乘积是目前改动最小的方案，能让权重随条件移动。

\subsection{固定工作点}

门控只有在阈值先于评测固定下来才有意义。候选阈值取自开发集的分数量点；当分位网格无法在方法之间对齐时，改用 $0.35$ 至 $0.97$、步长 $0.01$ 的加密扫描。每个候选都在校准来源上回放，工作点取"在校准失败比例不超过天花板的前提下恢复缺口最多"的那一个。我们同时报告松弛与严格两个天花板：只给一个数字会掩盖权衡的形状。阈值及其输入哈希在任何测试来源被读取之前就已写盘。

这个流程本身就是一次选择，我们按选择来对待它。一次扫描最多评估 63 个阈值，而用同一批既定了天花板、又从中挑出最优成员的数据，会把得到的风险压向乐观。Learn-then-Test \cite{angelopoulos2021learn} 正是为这类策略族上的选择做校准的，候选集覆盖 \cite{li2022pac} 与检测不确定性 \cite{su2024motcup} 也有相应保证；本文不采用它们。下文出现的每一个失败比例，都是对所指划分的观测值，而不是对未见来源的界。
